% BHU PYQ -- Manually transcribed & error-corrected
% Paper: BCH-221 : Business Regulatory Framework II
% Exam: B.Com. (Hons.) Semester IV Examination, 2022-23

\documentclass[11pt,a4paper]{article}
\usepackage[a4paper,top=2.5cm,bottom=2.5cm,left=2.8cm,right=2.8cm,headheight=14pt]{geometry}
\usepackage{amsmath,amssymb}
\usepackage[T1]{fontenc}
\usepackage[utf8]{inputenc}
\usepackage{lmodern,microtype,fancyhdr,enumitem,hyperref,lastpage,parskip}

\hypersetup{
    colorlinks=true,
    urlcolor=black,
    linkcolor=black,
    pdftitle={BCH-221: Business Regulatory Framework II -- BHU B.Com. (Hons.) Sem IV 2022-23}
}

\pagestyle{fancy}
\fancyhf{}
\renewcommand{\headrulewidth}{0.4pt}
\renewcommand{\footrulewidth}{0.4pt}
\fancyhead[L]{\small\textit{Banaras Hindu University}}
\fancyhead[R]{\small\textit{BCH-221: Business Regulatory Framework II}}
\fancyfoot[C]{\small Page \thepage\ of \pageref{LastPage}}

\newlist{parts}{enumerate}{2}
\setlist[parts,1]{label=(\alph*),leftmargin=*,itemsep=5pt,topsep=3pt}
\setlist[parts,2]{label=(\roman*),leftmargin=*,itemsep=3pt,topsep=2pt}
\newcommand{\pts}[1]{\hfill\textit{[#1]}}

\begin{document}

\begin{center}
  {\large\bfseries BANARAS HINDU UNIVERSITY}\\[2pt]
  {\normalsize B.Com. (Hons.) --- Semester IV Examination, 2022-23}\\[6pt]
  \rule{0.8\linewidth}{0.5pt}\\[6pt]
  {\large\bfseries Paper: BCH-221: Business Regulatory Framework II}\\[4pt]
  {\normalsize Time: 3 Hours\hfill Maximum Marks: 70}
\end{center}

\rule{\linewidth}{0.4pt}

\smallskip
\noindent\textit{Instructions: Answer all questions. All questions carry equal marks (14 marks each).}

\bigskip

\noindent\textbf{Question 1.} \\
Discuss various essential characteristics of a contract of sale under the Sale of Goods Act, 1930.\pts{14}

\centerline{\textbf{OR}}

Who is an Unpaid Seller? Explain the rights of an unpaid seller.\pts{14}

\medskip\rule{\linewidth}{0.2pt}

\noindent\textbf{Question-2.} \\
Explain the following:\pts{7+7}
\begin{parts}
  \item Registration of Partnership Firm
  \item Partnership Deed
\end{parts}

\centerline{\textbf{OR}}

Define Partnership. Explain different types of Partnership. Differentiate between Dissolution of Partnership and Dissolution of firm.\pts{14}

\medskip\rule{\linewidth}{0.2pt}

\noindent\textbf{Question 3.} \\
State the difference between ``Holder'' and ``Holder in due course''. Explain the privileges of a holder in due course under Negotiable Instruments Act, 1881.\pts{14}

\centerline{\textbf{OR}}

Explain Negotiable Instruments and write short notes on Bills of Exchange and Cheque.\pts{14}

\medskip\rule{\linewidth}{0.2pt}

\noindent\textbf{Question 4.} \\
What do you understand by Crossing of Cheque? Explain its significance. Explain various types of crossing.\pts{14}

\centerline{\textbf{OR}}

What is the difference between Discharge and Dishonour of Negotiable Instruments? Explain different types of dishonour of Negotiable Instruments.\pts{14}

\medskip\rule{\linewidth}{0.2pt}

\noindent\textbf{Question 5.} \\
Discuss the main provisions of Consumer Protection Act, 1986.\pts{14}

\centerline{\textbf{OR}}

Write short notes on the following:\pts{7+7}
\begin{parts}
  \item Sale and Agreement to Sell
  \item Consumer Protection Council
\end{parts}

\vfill
\rule{\linewidth}{0.4pt}
\begin{center}\small
  \href{https://bhu-exam.vercel.app/}{Click here to visit our website}\\[4pt]
  \href{https://me-aryan.vercel.app/}{Click here to visit our contributor}\\[4pt]
  \href{https://quashan-me.vercel.app/}{Click here to visit our contributor}
\end{center}

\end{document}
